\section{Formal evaluations and generic spectra}\label{sec:formal-generic}

We begin the proof of Theorem~\ref{thm:lifting} by constructing a
substitution $Z=(Z_1,Z_2)$ for the two free letters. Its coefficients will commute
with the old algebra $A$, so applying it to $(X,w)$ preserves their
polynomial identities. The matrix calculation focuses on $w$;
Sections~\ref{sec:division-realization} and~\ref{sec:generation}
realize the substitution in a division algebra of larger degree and
restore generation and both transverse directions for $X(Z)$.

Throughout the lifting argument, $k$ is a field of characteristic zero
containing all roots of unity, $A$ is a central division algebra of degree
$d$ over $k$, and
\[
 w\in A\langle\!\langle z_1,z_2\rangle\!\rangle.
\]
The letters commute with $A$, but not with each other.  Thus an element
is a family $(a_\omega)_\omega$ indexed by finite words, written
$\sum_\omega a_\omega z^\omega$, and multiplication uses concatenation
of words.  Fix an algebraic closure $K$ of $k$ and a splitting
$A\otimes_k K\simeq M_d(K)$; this uses the splitting and
scalar-extension properties of central simple algebras
\cite[Chapter~IV, Example~2.14(a) and Proposition~2.15]{MilneCFT}.

If $P=(P_1,P_2)$ is a pair of $n\times n$ indeterminate matrices, put
\[
 W_n(P)=\sum_\omega a_\omega\otimes P^\omega
       \in M_{dn}(K[[P]]).
\]
Here $K[[P]]$ is the ordinary commutative power series ring in the
$2n^2$ entries.  We omit the subscript $n$ when the size is clear.
The coefficient matrices act on the first tensor factor.  In
particular, every input arrow between blocks is a pair of scalar
matrices acting on the second factor.  It is not an arbitrary pair of
matrices on both factors.

In the spectral case we seek a substitution $Z=(Z_1,Z_2)$, with
constant term $R$, for which $W(Z)$ is conjugate to a split semisimple
constant output $W(R)$. If $\Lambda$ is the set of distinct eigenvalues
of $W(R)$, this gives
\[
 q(W(Z))=0,\qquad q(t)=\prod_{\lambda\in\Lambda}(t-\lambda).
\]
The division realization will identify $q$ with the minimal polynomial
over its central coefficient field. If $w$ already satisfies the
minimal polynomial of $w(0)$, substitution retains that relation and
the spectral correction is unnecessary.

For the spectral construction, we will find matrix sizes with a common
eigenvalue multiplicity, disjoint spectra at independent inputs, and
scalar spectral blocks under commuting Taylor perturbations. These
properties will make the eigenvalue responses of
Section~\ref{sec:responses} well-defined. We first specify the
evaluations and specializations used to establish them and to evaluate
triangular auxiliary inputs later.

\subsection{Convergence and specialization}\label{subsec:formal-conventions}

Let $E$ be a complete discretely valued field containing $K$, with
valuation $v$, such that $v(c)=0$ for $c\in K^\times$.  We use
$E=C((\varepsilon))$, or a finite extension of such a field, and
normalize $v(\varepsilon)=1$.  Matrix valuations mean the minimum of
the entry valuations.  A matrix is \emph{small} if this minimum is
positive.  A small generic point is
$P=\varepsilon P^\circ$, where all entries of $P^\circ$ are
algebraically independent over the previously specified constant
field.  Several independent points can use the same $\varepsilon$.

The finite extensions used to split spectra are separable in
characteristic zero. The valuation extends to each such extension,
which remains complete and discretely valued
\cite[Theorem~7.38 and Remark~7.43]{MilneANT}; with the normalization
$v(\varepsilon)=1$, its value group may be $\frac1e\mathbb Z$ for a
positive integer $e$.

\begin{lemma}[Admissible evaluations]\label{lem:admissible-evaluation}
Suppose a constant input pair $R$ is block upper triangular with $b$
blocks and small diagonal blocks.  Its strictly upper blocks may have
arbitrary entries in $E$.  Let $H$ be a matrix pair of formal series
of positive order in finitely many commuting or free variables.
Then $W(R+H)$ is defined coefficientwise by convergent sums in $E$.
The same assertion holds for products, finite determinants, and
derivatives of fixed order of finitely many such entry series.

These evaluations preserve formal entry identities.  They also
remain valid coefficientwise in an additional parameter $\eta$ if
the diagonal blocks have entries in $E[[\eta]]$ with small constant
terms, while the strictly upper blocks and the coefficients of $H$
have entries in $E((\eta))$.
\end{lemma}

\begin{proof}
Fix a coefficient of degree $r$ in the variables of $H$.  An input
path contributing to this coefficient uses at most $r$ positive-order
terms of $H$.  Between these insertions it follows the constant upper
triangular pair.  In each of the at most $r+1$ resulting intervals,
there are at most $b-1$ strictly upper steps.  Thus at most
$(r+1)(b-1)$ constant factors are strictly upper.  All remaining
constant factors come from small diagonal blocks.

Only finitely many coefficients of $H$ can occur at the chosen degree.
Their valuations, and those of the finitely many strictly upper
blocks, have a common lower bound.  If the original word has length
$L$, its contribution consequently has valuation at least
$cL-C_r$, where $c>0$ and $C_r$ is independent of the word.  There
are finitely many words of each length, so the sum converges.  The
same estimate applies after inserting the finitely many extra factors
needed for a product or a derivative of fixed order.  Completeness
allows regrouping of these sums, proving compatibility with
multiplication and with identities.

For the last assertion, first fix both the degree $r$ in $H$ and an
$\eta$-coefficient.  The bounded number of strictly upper steps and
of $H$-insertions bounds their total negative $\eta$-order.  Hence
only a bounded number of positive-order diagonal $\eta$-terms can
contribute, and only finitely many of their coefficients are used.
Every other factor is a small constant diagonal term.  The preceding
valuation estimate therefore applies again.  The resulting assertion
is coefficientwise; no interchange of the two Laurent topologies is
being asserted.
\end{proof}

We will use inverses of finite matrices as rational operations after
evaluation.  The following elementary distinction is important.
Although $(\varepsilon-Q)^{-1}$ has a formal expansion at $Q=0$ over
$E$, that expansion need not converge after $Q=\varepsilon Q^\circ$.
Accordingly, when transferring an identity involving an inverse, we
first multiply by its determinant.  We evaluate the resulting
numerator by Lemma~\ref{lem:admissible-evaluation}, and divide only
when the evaluated determinant is nonzero.  Formal inverse expansions
are used to define coefficients; they are not presumed to converge at
every subsequent specialization.

\begin{lemma}[Detection by independent residues]\label{lem:generic-detection}
Let
\[
 F(Q)=\sum_\alpha c_\alpha Q^\alpha\in E[[Q_1,\ldots,Q_N]],
 \qquad \inf_\alpha v(c_\alpha)>-\infty.
\]
Adjoin elements $Q_1^\circ,\ldots,Q_N^\circ$ whose residues are
algebraically independent over the residue field of $E$, and complete
with respect to the extended valuation.  Then
$F(\varepsilon Q^\circ)$ is defined, and is nonzero whenever $F$ is
nonzero.
\end{lemma}

\begin{proof}
The valuations $v(c_\alpha)+|\alpha|$ tend uniformly to infinity
with $|\alpha|$.  Their least value is therefore attained by a
nonempty finite set of monomials.  After scaling to this valuation,
the residue of their sum is a nonzero polynomial in the independent
residues of the $Q_j^\circ$.  This proves both assertions.
\end{proof}

In applications, the coefficients of a cleared numerator are finite
sums of products of base-point scalars and matrix-entry series.
The path estimate gives a uniform valuation bound in the auxiliary
coordinates.  Lemma~\ref{lem:generic-detection} therefore identifies
formal identities in these coordinates with identities at independent
small generic points.  The base-point valued field is constructed
first; the auxiliary residues are adjoined afterward.  This order
keeps their independence valid even after the base spectrum has been
split.  The argument can be applied separately at each Taylor degree.

We record two further field conventions.  Evaluation
\[
 K[[P]]\longrightarrow C(P^\circ)((\varepsilon)),
 \qquad P\longmapsto\varepsilon P^\circ,
\]
is injective, by homogeneous degree or by
Lemma~\ref{lem:generic-detection}.  It extends to the fraction field,
and a finite splitting field embeds into a finite algebraic extension
of the target.  Also, if $R$ is a small constant pair, then
\begin{equation}
K[[P]]\longrightarrow E[[\Delta]],\qquad P\longmapsto R+\Delta
 \label{eq:generic-taylor-embedding}
\end{equation}
is injective: reduction modulo the valuation ideal sends the image of
$f$ to $f(\Delta)$.  We use this observation for small block diagonal
$R$.  Denominators in the fraction field remain nonzero, so the map
extends to fraction fields.

Finally, the Taylor substitution $P\mapsto P+\eta v$ maps
$\operatorname{Frac}K[[P]]$ into
$\operatorname{Frac}K[[P]][[\eta]]$, with reduction the identity.
It extends to a finite separable extension $\Omega$ with values in
$\Omega[[\eta]]$: choose a primitive element and lift its defining
equation using its nonzero derivative.  The lift reducing to the
chosen primitive element is unique, by coefficient recursion.  This
convention lets us move a base point while retaining its labeled
eigenvalue branches.

\subsection{Generic division evaluations}\label{subsec:generic-division}

Generic matrix evaluations have a useful divisibility property for
their eigenvalue multiplicities.  We obtain it by first evaluating
inside a division algebra.

For $n\ge1$, choose independent central variables $\alpha,\beta$
and a primitive $n$th root of unity $\zeta_n\in k$.  Over
$k_0=k(\alpha,\beta)$ let $B_n$ be generated by $x,y$ with
\[
 x^n=\alpha,\qquad y^n=\beta,\qquad yx=\zeta_nxy.
\]
For $n=1$ this means $B_1=k_0$.  The reduced monomials
$x^iy^j$, $0\le i,j<n$, form a basis.  Indeed they span by the
relations, and after adjoining $n$th roots of $\alpha$ and $\beta$
their clock-and-shift matrices are linearly independent.  These
matrices span $M_n$, so $B_n$ is central simple of degree $n$.
This is the standard symbol construction
\cite[Corollary~2.5.5(1)]{GilleSzamuely2006}, with the inverse
root of unity for our convention.

\begin{lemma}\label{lem:symbol-division}
The central simple algebra $A\otimes_k B_n$ is division.  Its pure
rational scalar extensions and its central Laurent series extensions
are division as well.
\end{lemma}

\begin{proof}
Start with $A((x))$, where $x$ commutes with $A$, and let $\sigma$
fix $A$ and send $x$ to $\zeta_nx$.  The skew Laurent series ring
$A((x))((y;\sigma))$, with $ya=\sigma(a)y$, is a division ring by
the leading-coefficient construction of inverses.  The elements
$x^n,y^n$ are central and algebraically independent.  The defining
relations give a unital map of $A\otimes_k B_n$ into this division
ring.  Central simplicity makes the map injective.  A finite
dimensional domain over a field is division, since multiplication by
a nonzero element is an injective and hence surjective linear map.

For a central division algebra $D/F$, the rational extension $D(t)$
embeds in the division ring $D((t))$.  Iterating proves the rational
assertion.  Because $D$ has a finite $F$-basis,
$D\otimes_F F((\varepsilon))=D((\varepsilon))$, which is again
division by leading coefficients.
\end{proof}

Choose a basis $b_1,\ldots,b_{n^2}$ of $B_n$, and let
$c_{ij}$, $i=1,2$, $j=1,\ldots,n^2$, be independent variables.
Evaluate $w$ at
\[
 P_i=\varepsilon\sum_j c_{ij}b_j.
\]
This gives an element of a central division algebra of degree $dn$
over
\[
 F_n=k(\alpha,\beta,(c_{ij}))((\varepsilon)).
\]
After splitting $A$ and $B_n$ by constants independent of the $c_{ij}$,
the $b_j$ become a basis of $M_n$.  The change from the $c_{ij}$ to
the entries of $P_i/\varepsilon$ is therefore an invertible linear
change of independent coordinates.  The resulting evaluation is a
small generic scalar matrix evaluation in the sense fixed above.

Let $e_n$ denote the minimal-polynomial degree of this element.  Its
minimal polynomial is irreducible and separable over $F_n$.
Consequently $W(P)$ is semisimple over a splitting field, and every
distinct root has the same multiplicity
\begin{equation}
t_n=\frac{dn}{e_n}.
 \label{eq:generic-multiplicity}
\end{equation}
For completeness, if $D/F$ has degree $N$ and $a\in D$ has minimal
polynomial $q$ of degree $e$, then left multiplication by $a$ has
characteristic polynomial $q^{N^2/e}$: view $D$ as a vector space
over $F[a]$.  After splitting $D=M_N$, left multiplication acts by
the matrix $a$ on each of its $N$ columns.  Its characteristic
polynomial is therefore the $N$th power of the matrix characteristic
polynomial.  Comparing root multiplicities gives $N/e$ in the
matrix representation.  This also proves that $N/e$ is an integer.

Linear independence of $1,a,\ldots,a^{e-1}$ is preserved by field
extension.  Thus the same $e_n,t_n$ describe $W(P)$ over a splitting
field of $\operatorname{Frac}K[[P]]$.

\subsection{Nonconstant roots and stable multiplicity}
\label{subsec:stable-multiplicity}

Let $q_0$ be the minimal polynomial of $w(0)$ over $k$.  If
$q_0(w)=0$ as a free series, the subsequent spectral construction is
unnecessary; the unconstrained case of the lifting argument will use
this existing relation.  In the rest of this section assume
$q_0(w)\ne0$.

A nonzero homogeneous part of $q_0(w)$ is detected by scalar matrices
of every sufficiently large size.  To see this, choose a word of
length $L$ with nonzero coefficient and use a path of $L+1$ vertices,
placing its successive letters on the consecutive matrix-unit
arrows.  The $(1,L+1)$ entry detects precisely that coefficient.
Padding by zero blocks works for all larger sizes.  It follows that
$q_0(W(P))\ne0$ at every generic size $n\ge n_0$, for some $n_0$.

\begin{lemma}\label{lem:nonconstant-roots}
For $n\ge n_0$, every distinct generic eigenvalue of $W(P)$ has a
nonzero differential in the input entries.  Independent small generic
points of sizes at least $n_0$ have disjoint spectra.  Their spectra
also avoid the constant spectrum of $w(0)$.
\end{lemma}

\begin{proof}
Use the division evaluation over $F_n$.  Each derivation
$\partial/\partial c_{ij}$ acts coefficientwise on Laurent series
and extends uniquely to algebraic extensions in characteristic zero.
Uniqueness implies that it commutes with automorphisms of a splitting
field over $F_n$.  If all these derivations killed one root, they
would kill all its Galois conjugates.  Irreducibility makes these all
the distinct roots.  Thus every coefficient of the reduced
characteristic polynomial would have all its input derivatives zero.

The coefficient of $\varepsilon^r$ in such a characteristic
coefficient is a homogeneous polynomial of degree $r$ in the
$c_{ij}$.  In characteristic zero a polynomial whose partial
derivatives all vanish is constant.  All terms of positive
$\varepsilon$-degree therefore vanish.  The characteristic polynomial
is the characteristic polynomial of $w(0)$.  Semisimplicity then
implies $q_0(W(P))=0$, a contradiction.

Splitting $A,B_n$ uses constants independent of the inputs.  The
invertible basis change described above, together with the nonzero
factor $\varepsilon$, transports the differential assertion to
ordinary matrix-entry coordinates.  For independent points $P,Q$,
derivations in the $P$-coordinates kill the roots belonging to $Q$.
They cannot kill a root belonging to $P$.  No root can therefore
belong to both spectra.  The same argument excludes a constant root.
\end{proof}

Choose $p\ge n_0$ for which the positive integer $t_p$ is minimal
among all $t_n$, $n\ge n_0$, and put
\[
 t=t_p,\qquad \mathcal S=\{pN:N\ge1\}.
\]

\begin{lemma}\label{lem:stable-multiplicity}
For every $N\ge1$,
\[
 t_{pN}=t,\qquad e_{pN}=Ne_p.
\]
\end{lemma}

\begin{proof}
Specialize a size-$pN$ point to a diagonal sum of $N$ independent
size-$p$ points.  Lemma~\ref{lem:nonconstant-roots} gives $Ne_p$
distinct eigenvalues there.  The generic bound on minimal-polynomial
degree survives this specialization: it is expressed by the vanishing
of all $(e_{pN}+1)$-minors of the matrix whose columns are the
vectorizations of $1,W,\ldots,W^{e_{pN}}$.  These are formal entry
identities covered by Lemma~\ref{lem:admissible-evaluation}.
Consequently $e_{pN}\ge Ne_p$.  Minimality of $t$ gives the opposite
inequality,
$e_{pN}=dpN/t_{pN}\le dpN/t=Ne_p$.
\end{proof}

We henceforth use sizes in $\mathcal S$.  At size $n$, the minimal
polynomial always has degree at most $e_n=dn/t$ at every admissible
specialization.  The following consequence controls all commuting
Taylor deformations used below.

\begin{lemma}[Scalar branches at maximal root count]
\label{lem:scalar-branches}
Let $V(T)$ be a matrix over a commuting formal Taylor ring $E[[T]]$.
Suppose $V(0)$ is split semisimple with $e$ distinct eigenvalues and
the minimal-polynomial degree of $V(T)$ over
$\operatorname{Frac}E[[T]]$ is at most $e$.  Then a formal change of
basis reducing to the identity separates $V(T)$ into $e$ scalar
blocks.  The sizes of these blocks are their multiplicities at
$T=0$.
\end{lemma}

\begin{proof}
Use the eigenspace decomposition at zero.  Off-diagonal blocks can be
removed recursively in total Taylor degree: the equation for the
block between eigenvalues $\lambda_i,\lambda_j$ is multiplication
by $\lambda_i-\lambda_j$, a nonzero scalar.  We obtain blocks
$V_i(T)$ reducing to $\lambda_i I$.  Their characteristic polynomials
have pairwise unit resultants, because their reductions have different
roots.  Their minimal polynomials over the Taylor fraction field are
therefore pairwise coprime.  The minimal-polynomial degree of their
direct sum is the sum of their degrees.  There are $e$ nonempty
blocks and the total degree is at most $e$, so every block has
minimal-polynomial degree one.  Each $V_i(T)$ is scalar.  Its scalar
entry belongs to $E[[T]]$, since the block entries do.
\end{proof}

In particular, the roots at a small semisimple point with $e_n$
distinct roots have labeled scalar Taylor branches.  The splitting
field of the generic point embeds in their Taylor fraction field by
\eqref{eq:generic-taylor-embedding}; a labeling is chosen by reduction
to the roots at the base.

Let $\rho_n$ be the generic rank of the input Jacobian of
\[
 \operatorname{tr}W,\operatorname{tr}W^2,\ldots,
 \operatorname{tr}W^{dn}.
\]
Its minors are formal entry identities, so $\rho_n$ bounds the
corresponding specialized ranks.  At a maximal-root-count point the
scalar branches satisfy
\[
 d\operatorname{tr}W^j
   =jt\sum_{h=1}^{e_n}\lambda_h^{j-1}\,d\lambda_h.
\]
The first $e_n$ rows give an invertible Vandermonde matrix.  Hence
the trace Jacobian and the eigenvalue Jacobian have the same rank
there, including over a commuting Taylor fraction field.  We have
therefore obtained both ingredients needed for the response
calculations: scalar branches and specialization-stable rank bounds.
